\clearpage
\renewcommand{\sectioncolor}{quantum}
\section{``Apparent random'' mutation --- the word \emph{apparent} is doing real work}
\markboth{Mutation}{}

Your phrasing was exact. Mutation is \emph{apparently} random on two independent levels, and they
need separating because they sit in different tiers.

\subsection{Level 1 --- not uniform in space or time}
\noindent\Ttwo\;\emph{established; the mathematics exists and is simply not used here.}\medskip

\begin{center}\small
\begin{tabular}{@{}>{\raggedright\arraybackslash}p{8.4cm}>{\raggedright\arraybackslash}p{8.2cm}@{}}
\toprule
\rowcolor{quantum!12}
\textbf{\color{quantum}Finding} & \textbf{\color{quantum}Evidence}\\
\midrule
Mismatch repair efficiency tracks chromatin state, so mutation rate varies more than two-fold across
 the genome & Supek \& Lehner, \textit{Nature} \textbf{521}, 81 (2015)\\
$\sim$80+ distinct mutational signatures decompose the 96-channel substitution spectrum &
 Alexandrov et al., \textit{Nature} \textbf{500} (2013); PCAWG, \textit{Nature} \textbf{578} (2020)\\
CpG deamination produces $\sim$10-fold local hotspots & Lindahl, \textit{Nature} \textbf{362}, 709 (1993)\\
Mutation rate is \emph{regulated}: stress induces error-prone polymerases (SOS response) &
 Rosenberg and successors --- mutation rate is a controlled variable, not a constant\\
\rowcolor{alert!7}
Mutation bias correlates with gene essentiality in \textit{Arabidopsis} & Monroe et al.,
 \textit{Nature} \textbf{602} (2022) --- \textbf{contested}; treat as unsettled\\
\bottomrule
\end{tabular}
\end{center}

\begin{tier2}[title={The wrong object, and the right one}]
Standard practice models mutation as a \textbf{homogeneous Poisson process} with rate $\mu$. If the
rate varies across the genome and over time, the correct object is a \textbf{Cox process} --- a
Poisson process whose intensity is itself a random field:
\[
\lambda(x,t)\;=\;\lambda_0\,\exp\!\big(f(\text{chromatin},\,\text{replication timing},\,\text{transcription},\,\text{stress})\big)
\]
\begin{center}
\begin{tikzpicture}[font=\scriptsize]
% homogeneous
\node[anchor=east,font=\scriptsize\bfseries,text=slate] at (-0.3,1.5) {POISSON};
\draw[slate!50,line width=1pt] (0,1.5)--(15.2,1.5);
\foreach \x in {0.8,2.1,3.4,4.0,5.6,6.9,8.1,9.4,10.2,11.6,12.9,14.3}{
  \draw[slate,line width=1.1pt] (\x,1.33)--(\x,1.67);}
\node[anchor=west,font=\scriptsize,text=slate] at (15.4,1.5) {};
% cox
\node[anchor=east,font=\scriptsize\bfseries,text=measure] at (-0.3,0.0) {COX};
\draw[measure!35,line width=0.8pt,smooth,domain=0:15.2,samples=140]
  plot (\x,{0.45*sin(\x*38)+0.30*sin(\x*97+40)+0.55});
\draw[slate!50,line width=1pt] (0,0)--(15.2,0);
\foreach \x in {0.5,0.7,2.8,3.0,3.1,3.3,5.9,7.8,7.9,8.0,8.2,8.35,10.4,12.7,12.8,13.0,13.15,13.3,14.9}{
  \draw[measure,line width=1.1pt] (\x,-0.17)--(\x,0.17);}
\node[anchor=north,font=\scriptsize,text=measure,align=center] at (7.6,-0.42)
  {events cluster where the \emph{intensity field} is high --- hotspots, fragile sites, kataegis};
\end{tikzpicture}
\end{center}
\vspace{-2pt}
\textbf{Why this is not academic, for your AMR work:} if the intensity is a field rather than a
constant, the expected waiting time to a resistance mutation is \emph{not} $1/\mu$. Clustered
intensity makes rare-but-fast emergence far more likely than a Poisson model predicts. That changes
stewardship and dosing predictions, and the mathematics (log-Gaussian Cox processes, Hawkes processes
for clustered events) is completely standard and completely unused here.
\end{tier2}

\subsection{Level 2 --- the quantum substrate}
\noindent\Tthree\;\emph{this is where the mathematics genuinely runs out.}\medskip

This is where your question actually points, and it has a precise, citable literature.

\begin{evidence}[title={Löwdin's proposal, and its modern treatment}]
\textbf{Löwdin (1963)}, \textit{Rev.\ Mod.\ Phys.} \textbf{35}, 724, proposed that point mutations
arise from \textbf{proton tunnelling} across the hydrogen bonds of a base pair, producing rare
\textbf{tautomers} (G\*--C\*, A\*--T\*) that mispair when the replication fork arrives.

\textbf{Slocombe, Sacchi \& Al-Khalili}, \textit{Communications Physics} \textbf{5}, 109 (2022) ---
\emph{``An open quantum systems approach to proton tunnelling in DNA''} --- treated the G--C proton
transfer as a double well coupled to a thermal bath and found the tautomeric population reaches
$\approx 1.7\times10^{-4}$, equilibrating on \textbf{sub-picosecond} timescales: orders of magnitude
faster and more populated than a classical over-barrier estimate.
\end{evidence}

\begin{center}
\begin{tikzpicture}[font=\footnotesize]
% double well
\begin{scope}[shift={(0,0)}]
 \draw[slate!50,-{Stealth[length=4pt]}] (-3.3,0)--(3.5,0) node[right,font=\scriptsize]{proton position};
 \draw[slate!50,-{Stealth[length=4pt]}] (-3.3,0)--(-3.3,3.0) node[above,font=\scriptsize]{energy};
 \draw[quantum,line width=1.5pt,domain=-2.75:2.75,samples=180,smooth]
   plot (\x,{0.38*(\x*\x-2.6)*(\x*\x-2.6)/2.6 + 0.35});
 % levels
 \draw[bio,line width=1.2pt] (-2.15,0.62)--(-1.05,0.62);
 \node[font=\scriptsize,text=bio,anchor=east] at (-2.2,0.62) {canonical};
 \draw[alert,line width=1.2pt] (1.05,0.62)--(2.15,0.62);
 \node[font=\scriptsize,text=alert,anchor=west] at (2.2,0.62) {tautomer};
 % tunnelling arrow
 \draw[measure,line width=1.4pt,dashed,-{Stealth[length=5pt]}] (-1.15,0.62)--(1.15,0.62);
 \node[font=\scriptsize\bfseries,text=measure,anchor=south] at (0,0.70) {tunnelling};
 \node[font=\scriptsize,text=slate,anchor=north] at (0,2.55) {barrier $\gg k_BT$};
 \draw[slate,line width=0.6pt,{Stealth[length=3pt]}-{Stealth[length=3pt]}] (0,0.68)--(0,2.45);
 % bath
 \foreach \i in {0,...,9}{
   \draw[nano!60,line width=0.5pt] ({-2.7+\i*0.6},-0.45) sin ({-2.4+\i*0.6},-0.25) cos ({-2.1+\i*0.6},-0.45);}
 \node[font=\scriptsize,text=nano,anchor=north] at (0,-0.62)
   {the bath: ordered water, counterions, the helix itself --- \textbf{strongly coupled, structured, non-Markovian}};
\end{scope}
\end{tikzpicture}
\end{center}
\vspace{4pt}
\begin{center}\begin{minipage}{0.93\textwidth}\footnotesize\color{slate}
\textbf{Figure 4.}\; The Löwdin picture as an open quantum system. Everything hard is in the bath:
it is not weak, not memoryless, and not generic.
\end{minipage}\end{center}

\begin{tier3}[title={What mathematics is actually missing here}]
\begin{center}\small
\begin{tabular}{@{}>{\raggedright\arraybackslash}p{4.85cm}>{\raggedright\arraybackslash}p{5.85cm}>{\raggedright\arraybackslash}p{4.85cm}@{}}
\toprule
\rowcolor{alert!12}
\textbf{\color{alert}Missing object} & \textbf{\color{alert}Why it is hard} & \textbf{\color{alert}State of the art}\\
\midrule
Non-Markovian dynamics at \emph{strong} coupling &
 Lindblad needs weak coupling \emph{and} memorylessness --- both violated. HEOM is numerically exact
 but scales exponentially in bath complexity &
 \textbf{Genuinely open.} No tractable general theory\\
Decoherence in \emph{structured} environments &
 Tegmark-style estimates assume a generic bath. DNA's environment is ordered and correlated &
 No theory\\
\rowcolor{alert!6}
Quantum-origin error meeting a \emph{classical} proofreader &
 Kinetic proofreading (Hopfield 1974) is entirely classical. There is no theory of a classical
 corrector acting on a quantum-origin error &
 \textbf{Not even properly posed}\\
\bottomrule
\end{tabular}
\end{center}
\end{tier3}

\begin{critbox}[title={Be careful here --- and tell your students the same}]
The tunnelling is real physics. \textbf{The link to mutation is not closed.} The tautomer must
survive until the replication fork arrives \emph{and} the polymerase must misread it, and tautomer
lifetimes are short. Most spontaneous mutation is attributed to replication error, deamination and
oxidative damage, not tunnelling.

This is the same epistemic situation as the photosynthesis coherence claims in Part VII: genuine
quantum mechanics, unproven functional consequence. \textbf{Mark it as open, work on it if you want
to do mathematics, and do not cite it as established.}
\end{critbox}
